\section{SAT Encoding Correctness}
\label{app:sat-correctness}

This appendix records proof sketches for the two SAT encodings implemented in
the project. Let $G=(V,E)$ be a simple undirected graph with $n\geq 1$ vertices
and at least one edge. A labeling is a bijection
$f:V\rightarrow\{1,\ldots,n\}$, and its antibandwidth is
$B(f)=\min_{\{u,v\}\in E}|f(u)-f(v)|$.

\subsection{Pairwise encoding}
The assignment clauses give each vertex at least one label (ALO) and at most
one label (AMO). Therefore, every vertex receives exactly one label. The
permutation clauses apply the same ALO/AMO pattern to each label, so every
label is used by exactly one vertex. Consequently, the assignment variables
define a bijection.

Suppose a model satisfies the pairwise distance clauses at threshold $K$.
If an edge $\{u,v\}$ had endpoint labels $p$ and $q$ with
$|p-q|<K$, the clause $\neg x_{u,p}\vee\neg x_{v,q}$ would be present and
false. This is a contradiction. Conversely, a labeling whose edge distances
are all at least $K$ falsifies none of the distance clauses. Hence the
pairwise CNF is satisfiable exactly when a feasible labeling exists.

\subsection{Order encoding}
For each vertex, the chain clauses enforce
$o_{v,p}\Rightarrow o_{v,p+1}$. Therefore the true variables form a suffix.
The smallest true index is a unique label in $\{1,\ldots,n-1\}$; the all-false
chain represents label $n$. This establishes a one-to-one correspondence
between monotone chains and labels.

For a label $q$, the disjunction $N(v,q)$ in
Equation~\ref{eq:not-label} is true exactly when vertex $v$ does not receive
$q$. The permutation clause $N(u,q)\vee N(v,q)$ therefore excludes every
assignment in which two vertices receive the same label. With $n$ vertices
and $n$ labels, injectivity implies a bijection. The distance clauses use the
same disjunction to forbid every pair of endpoint labels whose difference is
less than $K$. Thus the order CNF is satisfiable exactly when a feasible
labeling exists.

\subsection{Faithfulness of AMO encodings}
For the pairwise AMO encoding, each clause
$\neg a_i\vee\neg a_j$ excludes exactly the assignments in which both
$a_i$ and $a_j$ are true. Thus the conjunction is satisfied if and only if
at most one input literal is true.

The sequential-counter AMO implementation is supplied by
\texttt{pysat.card.CardEnc.atmost} with bound one. Its auxiliary variables
represent whether a true input has appeared in a prefix. Monotonic propagation
of this prefix state prevents a second true input. Conversely, when at most
one input is true, setting each prefix state to the disjunction of the inputs
seen so far satisfies the counter clauses. The auxiliary variables therefore
do not change the feasible set projected onto the label variables.

\subsection{Equivalence theorem}
For both encodings and every $1\leq K<n$, the CNF is satisfiable if and only
if there is a bijective labeling $f$ satisfying $B(f)\geq K$. For $K\geq n$,
the problem is infeasible because the largest possible label difference is
$n-1$. The implementation handles this trivial case before building a CNF.

\subsection{Completeness of the multi-run search}
Let $B^*$ be the optimal antibandwidth. During the search, $\ell$ is the
largest verified objective found so far and $h$ is a threshold known to be
infeasible. The invariants are
\[
  \ell\leq B^*, \qquad h>B^*.
\]
They hold initially because $\ell=0$ and $h=n$. If $\ell<B^*$, the threshold
$K=\ell+1$ is feasible. A correct SAT/UNSAT response therefore cannot return
UNSAT at this threshold. A SAT response is decoded and verified, yielding a
value $B\geq K$; updating $\ell=B$ preserves the lower-bound invariant. An
UNSAT response at $K=\ell+1$ proves that no labeling reaches $K$, so
$B^*<K=\ell+1$. Combined with $\ell\leq B^*$, this gives
$B^*=\ell$ and the result is optimal.

After SAT at threshold $K$ with verified value $B$, all thresholds
$K<K'\leq B$ are feasible using the same labeling. Thus the next probe at
$B+1$ skips no unresolved threshold. Each iteration shrinks the interval
between $\ell$ and $h$ by at least one, so at most $n-1$ probes are required.
A timeout before any verified labeling produces TO; a timeout after an
incumbent produces FEAS with a verified lower bound and the best valid upper
bound known to the implementation.

\subsection{Trust boundary}
The verifier independently checks every reported labeling and recomputes its
objective. It does not independently certify UNSAT. Consequently, the
optimality claim depends on the SAT backend's UNSAT answer being correct;
the implementation's verification step protects the feasibility side of the
claim but is not an UNSAT proof checker.
